\addtocontents{toc}{\protect\Needspace{3\baselineskip}}
\originalchapter{横向振子如何确定整条弦}\label{ch:basic-lightcone}
对应 8.251 第 16 讲的经典解, 并衔接第 18 讲的正则量子化. 本章取自然单位, $0\leq\sigma\leq\pi$, 两端在所有空间方向均为 Neumann, $X^+=2\alpha'p^+\tau$, $p^+\ne0$.

\originalsection{模式与总动量的归一化}
反射边界把波动方程解组织成驻波,
\begin{equation}\label{eq:openmodes}
 X^I=x_0^I+2\alpha'p^I\tau
 +\ii\sqrt{2\alpha'}\sum_{n\ne0}\frac{\alpha_n^I}{n}
 \ee^{-\ii n\tau}\cos n\sigma,
 \qquad\alpha_0^I=\sqrt{2\alpha'}p^I.
\end{equation}
因 $\int_0^\pi\cos n\sigma\dd\sigma=0$, 积分 $P^{\tau I}=\dot X^I/(2\pi\alpha')$ 恰得总动量 $p^I$. 现实性要求 $\alpha_{-n}^I=(\alpha_n^I)^*$, 系数中的 $\ii/n$ 同时在 $n\mapsto-n$ 下改变符号, 故整个 $X^I$ 实. 将后续复振子当独立的两个实振幅时应保留这个关系.

计算导数给
\[
 \dot X^I\pm X'^I=\sqrt{2\alpha'}\sum_{n\in\mathbb Z}
 \alpha_n^I\ee^{-\ii n(\tau\pm\sigma)}.
\]
共形约束 $-2\dot X^+(\dot X^-\pm X'^-)+(\dot X^I\pm X'^I)^2=0$ 因而成为
\[
 \dot X^-\pm X'^-=\frac1{4\alpha'p^+}(\dot X^I\pm X'^I)^2.
\]
注意 $X^-$ 并不是额外一个独立的 Neumann 标量, 它由横向振幅的乘积决定.

\originalsection{Virasoro 模式从平方约束产生}
定义经典横向双线性
\[
 L_n^\perp=\frac12\sum_{k\in\mathbb Z}
 \alpha_{n-k}^I\alpha_k^I.
\]
将上述平方按 $\ee^{-\ii n(\tau\pm\sigma)}$ 展开并比较系数, 得
\[
 \sqrt{2\alpha'}\,\alpha_n^-=\frac{L_n^\perp}{p^+},
 \qquad 2\alpha'p^+p^-=L_0^\perp.
\]
零模是 $L_0^\perp=\alpha'p^Ip^I+\sum_{n>0}\alpha_{-n}^I\alpha_n^I$, 故经典不变质量
\[
 M^2=2p^+p^--p^Ip^I
 =\frac1{\alpha'}\sum_{n>0}\alpha_{-n}^I\alpha_n^I\geq0.
\]
在经典现实性条件下每项是模平方. 非零振动提供静止质量, 无振动解是退化的类光点状极限. 量子谱中出现的负截距尚未在这里产生, 不能把经典质量和量子质量层混用.

\originalsection{纵向解的局部可积性}
光锥约束直接给
\[
 \dot X^-=\frac{\dot X^I\dot X^I+X'^IX'^I}{4\alpha'p^+},
 \qquad X'^-=\frac{\dot X^IX'^I}{2\alpha'p^+}.
\]
对第二式取 $\tau$ 导数、对第一式取 $\sigma$ 导数, 两者之差正比于 $(\ddot X^I-X''^I)X'^I$, 因而为零. 所以可以在单连通区域沿任意路径积分求 $X^-$, 只需给一个位置常数 $x_0^-$. 闭弦的参数空间有非收缩圈, 局部可积不保证全局单值, 这将在闭弦章变成层级匹配.

\begin{example}
只有 $X^2=A\cos\tau\cos\sigma$ 非零且横向零动量时,
\[
 X^- =x_0^-+\frac{A^2}{8\alpha'p^+}\tau
 -\frac{A^2}{16\alpha'p^+}\sin2\tau\cos2\sigma.
\]
直接求导代入上面两式可验算. 纵向含二倍频, 正是横向场平方的结果. 零模给 $p^-=A^2/(16\alpha'^2p^+)$, 从而 $M^2=A^2/(8\alpha'^2)$; 也可从 $\alpha_{-1}\alpha_1$ 单模求得相同结果.
\end{example}

\originalsection{光锥坐标的正则结构}
在固定 $\tau$ 的约化相空间中, 横向场有
\[
 [X^I(\sigma),P^{\tau J}(\sigma')]=\ii\delta^{IJ}\delta_N(\sigma,\sigma'),
 \qquad[x_0^I,p^J]=\ii\delta^{IJ},\qquad[x_0^-,p^+]=-\ii.
\]
最后一个负号来自共轭动量 $p_-=-p^+$, 而非把 $p^+$ 当 $x^-$ 的普通正号动量. 在闭区间用满足 Neumann 条件的核
\[
 \delta_N(\sigma,\sigma')=\frac1\pi+\frac2\pi\sum_{n=1}^\infty
 \cos n\sigma\cos n\sigma'
\]
表示, 它在相应试验函数空间中是恒等核. 端点处的分布约定必须与偶延拓一致, 不能照搬整条直线的核而丢掉反射图像.

\originalsection{从场对易式得到振子对易式}
由 $P^{\tau J}=\dot X^J/(2\pi\alpha')$ 有
$[X^I,\dot X^J]=2\pi\alpha'\ii\delta^{IJ}\delta_N$.
将式\eqref{eq:openmodes}及其导数代入, 零模提供核的 $1/\pi$ 项. 对 $n,m>0$ 使用
$\int_0^\pi\cos n\sigma\cos m\sigma\dd\sigma=(\pi/2)\delta_{nm}$,
逐个投影振动部分, 得
\[
 [\alpha_m^I,\alpha_n^J]=m\delta^{IJ}\delta_{m+n,0}.
\]
以一个正频模式为例, $X_n=\ii\sqrt{2\alpha'}(\alpha_n\ee^{-\ii n\tau}-\alpha_{-n}\ee^{\ii n\tau})\cos n\sigma/n$,
$\dot X_n=\sqrt{2\alpha'}(\alpha_n\ee^{-\ii n\tau}+\alpha_{-n}\ee^{\ii n\tau})\cos n\sigma$. 若 $[\alpha_n,\alpha_{-n}]=n$, 其对易式恰为 $4\ii\alpha'\cos n\sigma\cos n\sigma'$, 与核展开逐项相等. 这也核对了 $\sqrt{2\alpha'}$ 和 $1/n$ 因子.

令 $a_n^I=\alpha_n^I/\sqrt n$、$a_n^{I\dagger}=\alpha_{-n}^I/\sqrt n$, 就得 $[a_n^I,a_m^{J\dagger}]=\delta_{nm}\delta^{IJ}$. 量子横向态具有正范数, 而光锥消除的时间及纵向模式不再被独立计入.
